\documentclass[11pt,a4paper,sans]{moderncv}
\moderncvstyle{classic}
\moderncvcolor{blue}
\usepackage[utf8]{inputenc}
\usepackage[scale=0.8]{geometry}

\name{Alexandro}{Disla}
\title{Lettre de Motivation}
\address{Rue Saint-Louis Jeanty \#14, Route de Frère}{Port-au-Prince, Haïti}{}
\phone[mobile]{(+509) 4148-3700}
\email{alexandrodisla@hotmail.com}
\extrainfo{GitHub: \url{https://github.com/AD0791}}

\begin{document}

\recipient{Équipe de Recrutement d'ACTED Haïti}{Port-au-Prince, Haïti}
\date{\today}
\opening{Madame, Monsieur,}
\closing{Je vous prie d'agréer, Madame, Monsieur, l'expression de mes salutations distinguées,}
\enclosure[Ci-joint]{Curriculum Vitae}

\makelettertitle

C'est avec un vif intérêt que je vous adresse ma candidature pour le poste d'Assistant(e) Base de Données au sein d'ACTED Haïti, sous la référence Réf : ASSISTBDD \_ 2607. Fort d'une expérience de plus de cinq ans faisant le pont entre le MEAL, l'administration de bases de données et le développement de logiciels, je dispose des compétences requises pour assurer la gestion et la maintenance des bases de données centrales d'ACTED tout en garantissant une qualité et une intégrité irréprochables des informations recueillies sur le terrain.

Mon parcours professionnel combine une solide base théorique en économie appliquée avec une expertise technique acquise dans la conception et l'optimisation d'infrastructures de données. Récemment, en tant qu'ingénieur logiciel chez Tekkod LLC, j'ai mené des projets complexes d'optimisation de performances de bases de données relationnelles. J'ai notamment résolu des boucles de requêtes N+1 inefficaces dans le projet « Steam The Streets », ce qui a permis de réduire le temps de latence de connexion de 36 secondes à moins de 0,5 seconde. De plus, j'ai mis en place des index composites pour éliminer les verrous transactionnels (locks et deadlocks) lors d'accès simultanés massifs et automatisé le nettoyage de données incohérentes. Ces compétences en profiling, en structuration de schémas de données et en traitement automatique de la qualité s'appliquent directement aux exigences d'ACTED en matière de gestion rigoureuse des bases de données et d'élimination systématique des doublons.

Par ailleurs, je maîtrise parfaitement les plateformes de collecte mobile KoboToolbox, ODK et CommCare. J'ai une longue pratique du codage avancé sous XLSForm, ce qui me permet d'intégrer des contraintes logiques strictes et des calculs d'indicateurs automatisés essentiels pour le ciblage des bénéficiaires et les enquêtes de suivi (telles que les enquêtes PDM ou Endline). Au cours de mes précédentes missions à la Fondation Caris et pour HANWASH, j'ai conçu des plans MEAL, administré des bases de données centrales (MySQL et mWater) et formé de nombreuses équipes d'enquêteurs de terrain, assurant une collecte à la fois opportune, précise et de qualité.

Ma parfaite maîtrise du français et du créole haïtien, associée à ma connaissance approfondie des réalités et contextes d'intervention locaux, me permettra de collaborer efficacement avec le Responsable MEAL d'ACTED, les équipes de terrain et de fournir un support technique fiable aux différents départements. Étant basé à Port-au-Prince, je suis immédiatement disponible et prêt à soutenir vos opérations.

Je serais honoré d'apporter mon expertise et ma rigueur méthodologique à ACTED afin de maximiser l'impact de vos programmes. En vous remerciant de l'attention accordée à ma candidature, je reste à votre entière disposition pour un entretien.

\makeletterclosing

\end{document}
